
\section{General rules}

\subsection{Game Master}
The Game Master (GM) guides the game. Even when the GM also plays a character, they may clarify or adjust a rule when necessary to keep the session moving and enjoyable. The GM resolves uncertainty, tracks points where applicable, and fills gaps that the written material does not cover.

The GM should balance two priorities:
\begin{itemize}
  \item The game should be fun.
  \item The game should be fair.
\end{itemize}

Give players information at the right time. Context their character would have known for a long time can be shared in advance, ideally at least a day before the session. Deliver recent, sensitive, or plot-critical information later, when it becomes relevant to play.

\subsection{Character}
Each player acts as a character during the game. A character sheet gives that player information and incentives, rather than a script that must be followed word for word.

\paragraph{Objectives}
Objectives are outcomes a character is rewarded for achieving. Some may be revealed only during the game; mark these clearly so the GM knows when to distribute them.

\paragraph{Secrets}
Secrets are facts that other players may discover. A secret can create opposing scoring incentives: the character may lose points if it is proven, while another player gains points for uncovering it. A character can also be rewarded for persuading others that a false claim is true. Resolve these using the group's final belief, not every theory expressed during play.

\paragraph{Actions}
Actions are behaviours or tasks the player should perform, or at least strongly consider, during the game. They give the roleplay texture and help move the story forward.

\paragraph{Background}
The background provides the character's history before the event. It may include narrative context, leads for the main mystery, information other players are seeking, and facts the character wants to conceal. For a complex story, organize it by day and time so its sequence is clear.

\subsection{How to play}
The group is trying to solve the main mystery, which should become clear enough to answer before the session ends. If the game defines a group loss condition, make that explicit to the players.

Stay in character as much as possible. Treat the background as real, follow your actions, and take your secrets seriously. Avoid asking out-of-character questions such as ``what is your secret?''; do not read your sheet aloud or reveal highly specific game information except when asking the GM for clarification. Use a private message or quiet conversation when the GM needs to resolve something secretly.

Points create incentives to complete secondary objectives, control what information you share, and investigate other characters. They should not replace the main mystery: if everyone protects too much information and the group fails to solve it, the game can end in a shared loss regardless of the points total.

When a secret is exposed, the affected character has less reason to hide that fact and may become more proactive. Players can reveal useful information while protecting themselves with a plausible alternative explanation. Clues may also be hidden in the physical spaces, so players should be encouraged to search and interact with the setting.

\paragraph{Suggested game phases}
\begin{enumerate}
  \item Introductions, in which everyone establishes their character.
  \item A first period of social play, relationships, and low-stakes hints.
  \item An inciting event that makes the main mystery unmistakable.
  \item Investigation, private consultations with the GM when needed, and a final vote or accusation.
\end{enumerate}

\subsection{How to create a game}
Start with the answer to the main mystery: who committed the act, what happened, or what the supernatural or hidden truth really is. Build the story backwards from that answer. Connect every character to at least one other character, and give at least half of them information that can contribute to solving the main mystery.

Write the complete background before distributing character-specific information. Include a dated timeline and err on the side of giving the GM too much context. Then tailor the relevant fragments into the character sheets, making sure each player has enough information to make meaningful choices.

Mark leads, contradictions, and possible clues in the complete background. Add red herrings, misleading accounts, and smaller twists deliberately, but ensure the true solution remains reachable through evidence. Decide each team's victory and loss conditions in advance, and reveal them to players at the appropriate moment.
